\documentclass[../main.tex]{subfiles}
\begin{document}

%========================================================================================
% FILENAME: 04_invariants.tex (v1.0.0)
%
% §4 Floor Invariants — burn strictly raises phi; redemption, issuance, maturity
% conversion preserve it; the floor is globally non-decreasing.  The single
% order-sensitive cycle input is isolated in prop:invariants:burn-order-residual.
%
% Proof budget: one body proof (floor-non-decrease, an assembly); the four supporting
% props carry \proofsketch margins only.  Silences applied; cycle only.  First-use index
% qualifier carried in burn-order-residual.
%
% Change history lives in version control (see the versioning policy in main.tex).
%========================================================================================

\section{Floor Invariants}
\label{sec:attestation:invariants}

Burns strictly raise the floor; redemption, issuance, and maturity conversion preserve it
exactly; the floor is therefore globally non-decreasing
(\zcref{thm:invariants:floor-non-decrease}). Among all cycle inputs, only the attestation
boundary state is order-sensitive within a cycle, and that dependence is resolved in
\zcref{prop:invariants:burn-order-residual}.

\begin{proposition}[Burn Increases Floor]
    \label{prop:invariants:burn-increases-floor}
    \proofsketch{$\reservepool{}$ unchanged; the denominator decreases.}%
    After a burn of $x > 0$ \liveclass{} units (\zcref{def:operations:burn}),
    $\ratefloor{}^{+} = \reservepool{}/(\totsupply{} - x) > \ratefloor{}$.
\end{proposition}

\begin{proposition}[Burn-Order Residual]
    \label{prop:invariants:burn-order-residual}
    Among all cycle inputs, only $\set{\attest{a}}$ depends on the intra-cycle sequence of
    operations, via the strictly increasing floor
    (\zcref{prop:invariants:burn-increases-floor}). The index $a$ is an opaque address
    (\zcref{def:interface:address-space}); the binding of addresses to any identity is the
    importing system's import-time interpretation. Two reductions remove the dependence:
    \begin{enumerate}[nosep]
        \item \textbf{Boundary-state provision (mechanism-agnostic).} The operator commits
            $\set{\attest{a}^{(k)}}$ as part of the cycle-$k$ settlement, auditable against
            floor non-decrease (\zcref{thm:invariants:floor-non-decrease}) and the supply
            identity (\zcref{def:model:supply}); the intra-cycle sequencing is out of this
            layer's scope. Continuous immediacy is retained.
        \item \textbf{Settlement-pinned clearing (order-free).} Every burn settles at the
            floor committed by the most recent settlement event
            (\zcref{postc:interface:monotonicity}), making $\attest{a}$ a commutative
            counter between events. With settlement events at cycle boundaries this is
            per-cycle batch clearing at a single frozen $\ratefloor{}^{(k)}$, and the
            within-cycle timing bound $1/(1-\cycleburnfrac{k})$
            (\zcref{sec:verification:timing}) becomes exact per batch; a sparser event
            family widens the lag and is conservative under
            \zcref{prop:interface:conservative-valuation}.
    \end{enumerate}
    Under either reduction every cycle quantity is a deterministic function of the boundary
    state.
\end{proposition}

\begin{proposition}[Redemption Preserves Floor]
    \label{prop:invariants:redemption-preserves-floor}
    \proofsketch{$(\reservepool{}-x\ratefloor{})/(\totsupply{}-x)=\ratefloor{}$.}%
    After redemption of $x > 0$ \liveclass{} units (\zcref{def:operations:redemption}),
    $\ratefloor{}^{+} = \ratefloor{}$.
\end{proposition}

\begin{proposition}[Issuance Preserves Floor]
    \label{prop:invariants:issuance-preserves-floor}
    \proofsketch{$\Delta\totsupply{} = \Delta\reservepool{}/\ratefloor{}^{-}$, so the ratio
    is unchanged; the split is internal to $\Delta\totsupply{}$ and does not enter it.}%
    After cycle minting with escrow $\Delta\reservepool{}$
    (\zcref{alg:operations:cycle-processing}),
    $\ratefloor{}^{+} = \ratefloor{}^{-}$.
\end{proposition}

\begin{proposition}[Maturity Conversion Preserves Floor]
    \label{prop:invariants:maturity-preserves-floor}
    \proofsketch{Maturity conversion relabels within $\totsupply{}$ and leaves
    $\reservepool{}$ untouched (\zcref{def:maturity:conversion}).}%
    After maturity conversion, $\ratefloor{}^{+} = \ratefloor{}$.
\end{proposition}

\begin{theorem}[Floor Non-Decrease]
    \label{thm:invariants:floor-non-decrease}
    $\ratefloor{}^{(0)} \leq \ratefloor{}^{(1)} \leq \cdots$, and within each cycle
    $\ratefloor{t} \geq \ratefloor{}^{(k)}$ for all $t$ in cycle $k$.
\end{theorem}
\begin{proof}
    Within a cycle no minting occurs: burns strictly increase $\ratefloor{}$
    (\zcref{prop:invariants:burn-increases-floor}) and redemptions preserve it
    (\zcref{prop:invariants:redemption-preserves-floor}). Across boundaries, issuance
    preserves it (\zcref{prop:invariants:issuance-preserves-floor}) and maturity conversion
    preserves it (\zcref{prop:invariants:maturity-preserves-floor}).
\end{proof}

%========================================================================================
% END OF FILE: 04_invariants.tex (v1.0.0)
%========================================================================================

\end{document}
